\subsection{FINITENESS CRITERION}

We retain the notation of the preceding nos., and assume that S is finite.

THEOREM 2. The following properties are equivalent:

\begin{enumerate}
\def\labelenumi{(\arabic{enumi})}
\item
  W is finite.
\item
  The form \(B_{M}\) is positive and non-degenerate.
\item
  \(\Longrightarrow\) (2). Let \(S = \bigcup_{i} S_{i}\) be the decomposition of \(S\) into connected components (Chap. IV, § 1, no. 9), and let \(W = \prod_{i} W_{i}\) be the corresponding decomposition of \(W\) . The space \(E\) can be identified with the direct sum of the spaces \(E_{i} = R^{S_{i}}\) , and \(B_{M}\) can be identified with the direct sum of the corresponding forms \(B_{M_{i}}\) . We are thus reduced to the case when \((W, S)\) is irreducible. Since \(W\) is assumed to be finite, \(E\) is a semi-simple \(W\) -module (Appendix, Prop. 2). By the Cor. to Prop. 5, it follows that \(E\) is absolutely simple. Let \(B'\) be a positive non-degenerate form on \(E\) , and let \(B''\) be the sum of its transforms under \(W\) . Since \(B''\) is invariant under \(W\) , it is proportional to \(B_{M}\) ( \(\S 2\) , no. 1, Prop. 1); since \(B_{M}(e_{s}, e_{s}) = 1\) for all \(s \in S\) , the coefficient of proportionality is \(>0\) , and since \(B''\) is positive so is \(B_{M}\) , which proves (2).
\item
  \(\Longrightarrow\) (1). If \(B_{M}\) is positive non-degenerate, the orthogonal group \(\mathbf{O}(B_{M})\) is compact (Integration, Chap. VII, §3, no. 1). Since \(\sigma(W)\) is a discrete subgroup of \(\mathbf{O}(B_{M})\) (Cor. 3 of Th. 1), it follows that \(\sigma(W)\) is finite, hence so is W. Q.E.D.
\end{enumerate}

The following result has been established in the course of the proof:

COROLLARY. If \((\mathrm{W},\mathrm{S})\) is irreducible and finite, \(\mathbf{E}\) is an absolutely simple W-module.

The finiteness criterion provided by Th. 2 permits the classification of all finite Coxeter groups (cf.~Chap. VI, §4). We restrict ourselves here to the following preliminary result:

PROPOSITION 8. If W is finite, the graph of (W, S) is a forest (Chap. IV, Appendix).

If not, this graph would contain a circuit \((s_{1},\ldots,s_{n})\) , \(n\geqslant3\) . Putting \(m_{i}=m(s_{1},s_{i+1})\) , \(1\leqslant i<n\) , and \(m_{n}=m(s_{n},s_{1})\) , this means that \(m_{i}\geqslant3\) for all i. Let

\[
x = e _ {s _ {1}} + \dots + e _ {s _ {n}}.
\]

Then \(\mathrm{B_M}(x,x) = n + 2\sum_{i < j}\mathrm{B_M}(e_{s_i},e_{s_j})\) . Now

\[
\mathrm{B} _ {\mathrm{M}} (e _ {s _ {i}}, e _ {s _ {i + 1}}) = - \cos \frac {\pi}{m _ {i}} \leqslant - \cos \frac {\pi}{3} \leqslant - \frac {1}{2},
\]

and similarly for \(\mathrm{B}_{\mathbb{M}}(e_{s_n}, e_{s_i})\) . Since the other terms in the sum are \(\leqslant 0\) , we obtain

\[
\mathrm{B} _ {\mathrm{M}} (x, x) \leqslant n - n = 0,
\]

contrary to the fact that \(B_{M}\) is positive non-degenerate.

COROLLARY. If (W, S) is irreducible and finite, its graph is a tree.

Indeed, a connected forest is a tree.

Comparison with the results of §3.

First of all let \((W,S)\) be a finite Coxeter group. Denote by \((x|y)\) the form \(\mathrm{B}_{\mathbf{M}}(x,y)\) ; by Th. 2, this is a scalar product on E. For all \(s \in S\) , let \(H_{s}\) be the hyperplane associated to the orthogonal reflection \(\sigma_{s}\) , and let \(\mathfrak{H}\) be the family of hyperplanes \(w(\mathrm{H}_{s})\) , for \(s \in S\) , \(w \in W\) . Let \(C_{0}\) be the set of \(x \in E\) such that \((x|e_{s}) > 0\) for all \(s \in S\) . Finally, identify W (by means of \(\sigma\) ) with a subgroup of the orthogonal group \(\mathbf{O}(\mathbf{E})\) of the space E.

PROPOSITION 9. With the preceding notations, W is the subgroup of O(E) generated by the reflections with respect to the hyperplanes of \(\mathfrak{H}\) . It is an essential group (§3, no. 7) and C0 is a chamber of E relative to \(\mathfrak{H}\) .

The first assertion is trivial. On the other hand, if \(x \in E\) is invariant under W, it is orthogonal to all the \(e_{s}\) , and hence is zero; this shows that W is essential. Finally, the isomorphism \(E \to E^{*}\) defined by \(B_{M}\) transforms \(C_{0}\) to the set C of no. 4; the property \((P_{n})\) proved there shows that, for all \(w \in W\) and all \(s \in S\) , \(w(C_0)\) does not meet \(H_s\) . It follows that \(C_0\) is contained in the complement U of the union of the hyperplanes of \(\mathfrak{H}\) , and since \(C_0\) is connected, open and closed in U, it is a chamber of E relative to \(\mathfrak{H}\) . Q.E.D.

We can apply to W and \(C_{0}\) all the properties proved in §3. In particular, \(\overline{C_{0}}\) is a fundamental domain for the action of W on E (in other words, the cone U defined in no. 6 is equal to the whole of E).

Conversely, let E be a finite dimensional real vector space, provided with a scalar product \((x|y)\) and let W be an essential finite group of displacements of E leaving 0 fixed; assume that W is generated by reflections. Let \(C_{0}\) be a chamber of E with respect to W (cf.~§3), and let S be the set of orthogonal reflections relative to the walls of \(C_{0}\) . Then (W, S) is a finite Coxeter system (§3, no. 2, Th. 1). Moreover, if \(s \in S\) , denote by \(H_{s}\) the wall of \(C_{0}\) corresponding to s, and denote by \(e_{s}\) the unit vector orthogonal to \(H_{s}\) and on the same side of \(H_{s}\) as \(C_{0}\) . If \((m(s, s'))\) denotes the Coxeter matrix of (W, S), Props. 3 and 7 of §3 show that

\[
(e _ {s} | e _ {s ^ {\prime}}) = - \cos (\pi / m (s, s ^ {\prime}))
\]

and that the \(e_{s}\) form a basis of E. The natural representation of W on E can thus be identified with the representation \(\sigma\) of no. 3.

